\chapter{Úvod}
\label{ch:uvod}

Udržování zdravého životního stylu a optimální tělesné kondice představuje klíčový aspekt moderního pojetí kvality života.
Tento význam je zdůrazněn nejen individuálními cíli, jako je zvýšení energie či osobní spokojenosti, ale také aktuálními zdravotními trendy populace.
Podle dostupných dat se Češi umisťují na čtvrtém místě v žebříčku evropských zemí z hlediska výskytu nadváhy a obezity.
Celkem 55,4~\% české populace má tělesnou hmotnost vyšší než doporučené normy.
Horší jsou pouze Řecko, Chorvatsko a Malta, kde tento podíl dosahuje až 59,6~\%~\cite{evropskaMapaObezity2019}.

Motivace k osvojování udržitelných stravovacích a pohybových návyků je proto více než aktuální a různorodá, a to od prevence civilizačních chorob až po snahu o zlepšení celkové fyzické kondice a životní pohody.
Dosažení těchto cílů však vyžaduje systematický a informovaný přístup, který přesahuje pouhé počáteční odhodlání a opírá se o vědecké poznatky o výživě, energetickém příjmu a výdeji, stejně jako o praktické nástroje pro monitorování a optimalizaci životního stylu.

Základním principem regulace tělesné hmotnosti je energetická bilance.
Tento koncept definuje vztah mezi kaloriemi přijatými ve stravě a energií, kterou tělo spotřebuje na svůj provoz a veškerou fyzickou aktivitu.
Ačkoliv je tato myšlenka ve své podstatě jednoduchá, její praktická aplikace naráží na řadu překážek.
Mnoho lidí své úsilí vzdává, protože manuální zaznamenávání stravy je vnímáno jako zdlouhavé a nepraktické pro každodenní použití.
Bez systematického přehledu se pak snadno uchylují k neefektivním dietním zkratkám, které nevedou k dlouhodobým výsledkům.

Moderní technologie nabízejí jedno z možných řešení těchto problémů.
Zejména mobilní aplikace se staly rozšířeným nástrojem, který dokáže uživatele provázet jejich snahou, poskytovat jim okamžitou zpětnou vazbu a uchovávat data potřebná pro kvalifikovaná rozhodnutí o jídelníčku a pohybové aktivitě.
Současně se rozvíjí oblast umělé inteligence a počítačového vidění, která umožňuje částečně automatizovat dosud manuální úlohy, například odhad složení jídla z fotografie.

Náplní této práce je prozkoumat potenciál nejnovějších pokroků v oblasti umělé inteligence pro zjednodušení sledování kalorického příjmu.
Hlavním cílem je navrhnout a částečně realizovat aplikaci, která využívá schopností pokročilých modelů umělé inteligence k automatické analýze jídla z fotografie, čímž odstraňuje část bariér spojených s ručním zapisováním a činí proces správy jídelníčku efektivnější a uživatelsky přijatelnější.

Dílčími cíli práce jsou provedení rešerše současné odborné literatury v oblasti digitálních nástrojů pro výživu a využití umělé inteligence v této doméně, dále analýza stávajících řešení na trhu, která umožní identifikovat silné a slabé stránky konkurence a definovat prostor pro inovaci.
Na tuto analýzu navazuje výběr vhodných technologií a nástrojů, které budou pro potřeby implementace dostatečně výkonné a zároveň prakticky použitelné.
Neméně důležitou fází práce je návrh softwarové architektury pro zajištění spolehlivého chodu a efektivní komunikace mezi jednotlivými komponentami systému a vytvoření přehledného a intuitivního uživatelského rozhraní.

Na úvodní kapitolu navazuje část věnovaná rešerši odborných zdrojů a existujících mobilních aplikací pro sledování stravy.
Ta vytváří kontext pro návrh vlastního řešení a umožňuje zasadit diskutovanou aplikaci do širšího rámce digitálního zdraví a výzkumu v oblasti výživy.
